\chapter{Collection: register and capture}

\eyebrow{acquisition with verifiable pairing}

\vspace{4pt}
During collection, a carcass is registered before images are captured. The image
inherits the pairing from the carcass, not from a filename.

\section{Create a batch}

A \textbf{batch} is one collection session (a slaughter day, a site). It is
created from the Project dashboard or through an import. Each batch groups the
carcasses collected together.

\section{Register a carcass}

The batch is selected and its \textbf{Overview} view is opened. A carcass is
registered with the following fields:

\begin{center}
\begin{tabularx}{\linewidth}{@{}l X@{}}
\toprule
\textbf{Field} & \textbf{Meaning} \\
\midrule
Physical tag \textbf{(required)} & The readable ID on the physical label. It is
      unique per batch and constrains against the pairing error. \\
Animal (lab) & The animal ID in the laboratory spreadsheet, the explicit link. \\
Treatment & Diet / treatment group. \\
Stratum & Sampling stratum, for stratified coverage. \\
\bottomrule
\end{tabularx}
\end{center}

\begin{caution}[title=The physical tag is mandatory]
A carcass cannot be created without a physical tag, and the same tag cannot be
used twice in one batch. This makes the image-to-animal link verifiable and
prevents the duplicate-label failure of the earlier collection.
\end{caution}

\section{Capture images}

The carcass is selected and its \textbf{Images} view is opened
(Figure~\ref{fig:carcass-images}). The camera is started, a view angle is chosen
(posterior, lateral, dorsal), a frame is captured, the preview is reviewed, and
the image is saved. The saved image is written to the carcass's folder with its
\tele{carcass\_id} already set.

\shot{carcass-images.jpeg}{The Images view of a carcass. Left: webcam capture with
a view selector and a capture preview; the \tele{Save (pair to \#IMG\_6250)} button
makes the pairing explicit. Right: for the selected carcass, the Inspector shows
the full record (identification, physical reference, model analysis and
conformation).}{carcass-images}

The Inspector on the right (Figure~\ref{fig:carcass-images}) holds the carcass's
complete record. Once the image has been analysed, it also shows the model
results: \tele{Fat (model)}, whether the background was removed, the experimental
grade, and the conformation convexities.

\begin{note}[title={Kinect (RGB + depth), optional}]
If a Kinect and its libraries are present, the capture panel offers an
\tele{RGB+depth} capture that also records a depth map. Without it, the panel
falls back to webcam capture and shows \tele{Kinect unavailable (webcam only)}.
\end{note}

\section{Edit carcass data later}

The carcass \textbf{Overview} view is an editable form for every field, including
the physical reference measured on the animal (fat thickness, GR, loin eye area).
These are entered when the laboratory measurements become available; they are
required by the physical-correlation analysis (Station~7).
